\documentclass{report}
\usepackage{amsmath,amssymb}
\setlength{\parindent}{0mm}
\setlength{\parskip}{1em}
\begin{document}
\begin{center}
\rule{6in}{1pt} \
{\large Andrea Franceschini \\
{\bf Multilevel Approaches for FSAI Preconditioning}}

Department of Civil and Environmental Engineering \\ University of Padova \\ via Trieste 63 \\ 35121 Padova (Italy)
\\
{\tt andrea.franceschini@dicea.unipd.it}\\
Victor, A. Magri\\
Carlo Janna\\
Massimiliano Ferronato\end{center}

The solution of a large size sparse linear system of equations
\begin{equation}
\mathbf{A} x = b,
\end{equation}
where $\mathbf{A}$ is a SPD matrix, requires the use of iterative methods based on Krylov
subspaces. It is well known, however, that the use of these methods alone generally
provide poor convergence. This can be remedied by using their preconditioned form, which
requires the definition of a preconditioner, i.e., a linear transformation that
approximates the action of $\mathbf{A}^{-1}$.

In this paper, two new preconditioners built upon the dynamic FSAI algorithm
[Janna et al., 2015] are considered: the BTFSAI - Block Tridiagonal Factorized Sparse
Approximate Inverse - and the DDFSAI - Domain Decomposition Factorized Sparse Approximate
Inverse. The first one is based on the successive application of the FSAI
algorithm on the
matrix reordered in a block tridiagonal structure and the second one is based on the
combination of a two level domain decomposition approach and the FSAI preconditioner,
which is used for computing the approximate inverse of the inner
submatrices and the Schur
complement.

In the Block Tridiagonal preconditioner - BTFSAI - the matrix $\mathbf{A}$ is first
reordered by the reverse Cuthill-McKee algorithm, with the aim to reduce its bandwidth,
and then it is divided into a block tridiagonal structure according to a given number
of blocks. After that, the block LDU decomposition of this new matrix is calculated
where the inverses of the diagonal blocks are approximated explicitly by the adaptive
FSAI algorithm.

With this approach, the global matrix is subdivided into small blocks, that should be
approximated with less effort. Though this method is intrinsically sequential,
parallelism can be exploited for each block by the use of the FSAI algorithm.
Calling $\mathbf{A}_i$ the SPD diagonal submatrix and $\mathbf{B}_i$ the upper diagonal
submatrix of the $i-th$ block, the block LDU decomposition reads

\footnotesize
\begin{equation}
\mathbf{A} =
\begin{bmatrix}
\mathbf{I}_1 & & & \\
\mathbf{B}^T_1 \mathbf{S}^{-1}_1 & \mathbf{I}_2 & & \\
& \ddots & \ddots & \\
& & \mathbf{B}^T_{n-1} \mathbf{S}^{-1}_{n-1} & \mathbf{I}_{n}
\end{bmatrix}
\begin{bmatrix}
\mathbf{S}_1 & & & \\
& \mathbf{S}_2 & & \\
& & \ddots & \\
& & & \mathbf{S}_n
\end{bmatrix}
\begin{bmatrix}
\mathbf{I}_1 & \mathbf{S}^{-1}_1 \mathbf{B}_1 & & \\
& \mathbf{I}_2 & \mathbf{S}^{-1}_2 \mathbf{B}_2 & \\
& & \ddots & \ddots \\
& & & \mathbf{I}_n
\end{bmatrix}
\end{equation}
\normalsize
where
\begin{equation}
\mathbf{S}_i = \mathbf{A}_i - \mathbf{B}^T_{i-1} \mathbf{S}^{-1}_{i-1}
\mathbf{B}_{i-1}, \qquad \qquad i = 2
\dots n,
\end{equation}
is the $i-th$ block Schur complement and $\mathbf{S}_1 = \mathbf{A}_1$,
by definition. Let
$\mathbf{S}^{-1}_i = \mathbf{F}^T_i \mathbf{F}_i$ be the approximate
inverse of the Schur complement
computed by the adaptive FSAI algorithm and let the matrix $\mathbf{H}_i$ be equal to the
product $\mathbf{F}_i \mathbf{B}_i$, it follows that the Schur complement recursively as
$\mathbf{S}_i = \mathbf{A}_i - \mathbf{H}^T_{i-1} \mathbf{H}_{i-1}$.
Using these relations,
the approximate inverse of the matrix $\mathbf{A}$ reads
\footnotesize
\begin{equation}
\mathbf{A}^{-1} \approx \mathbf{M}^{-1} =
\begin{bmatrix}
\mathbf{F}^T_1 & - \mathbf{F}^T_1 \mathbf{H}_1 \mathbf{F}^T_2 & & \\
& \mathbf{F}^T_2 & - \mathbf{F}^T_2 \mathbf{H}_2 \mathbf{F}^T_3 & \\
& & \ddots & \ddots \\
& & & \mathbf{F}^T_n
\end{bmatrix}
\begin{bmatrix}
\mathbf{F}_1 & & & \\
- \mathbf{F}_2 \mathbf{H}^T_1 \mathbf{F}_1 & \mathbf{F}_2 & & \\
& \ddots & \ddots & \\
& & -\mathbf{F}_n \mathbf{H}^T_{n-1} \mathbf{F}_{n-1} & \mathbf{F}_n \\
\end{bmatrix}.
\end{equation}
\normalsize
The main drawback of this approach is the need of computing $n - 1$ Schur complements
along with their FSAI approximations. While theoretically these matrices should be SPD,
the use of approximate inverses may cause the loss of this property. Indeed, for some
matrices, we can have indefinite Schur complements if the approximate inverses are not
accurate enough. In this case, the construction of the preconditioner should be restarted
allowing for a larger number of nonzero entries per row.

Roughly speaking, domain decomposition refers to a number of techniques which solve a
problem defined on a large domain $\Omega$ via the solution of smaller problems defined
over subdomains $\Omega_k, k = 1,\ldots,n$. Each problem can be solved independently and
the solutions can be gathered to build the outcome of the original problem.

In the domain decomposition FSAI preconditioner - DDFSAI - the graph of the original
matrix is reordered according to the k-way partition algorithm implemented in the METIS
software library. With this method, the number of edges connecting a subdomain to the
others is minimized, thus reducing the pieces of information that need to be communicated
among different blocks. At the same time, the size of subdomains is kept to be
approximately of the same order in favour of providing a good load
balance among different
threads, which are going to work in one or possibly more blocks. After this, each
independent block of the matrix is reordered according to the reverse Cuthill-McKee
algorithm in order to reduce its bandwidth.

The algorithm for constructing the DDFSAI preconditioner is almost the same as BTFSAI
restricted to two blocks, with the only difference consisting in the initial ordering of
the matrix, which in the current case uses the METIS library besides of the reverse
Cuthill-McKee ordering.

The results show that the BTFSAI technique is very promising in terms of decreasing
the total number of iterations and time for achieving the solution. Further, with
the DDFSAI technique, the computational pain is smaller, but stability and scalability
are generally better. We remark that these strategies are under development and
further computational improvements are being addressed, for instance a hybrid
MPI/OpenMP implementation.

\section*{References}
M. Benzi. Preconditioning techniques for large linear systems: A survey. {\em Journal
of Computational Physics}, 182(2):418 { 477, 2002.

C. Janna and M. Ferronato. Adaptive pattern research for block FSAI preconditioning.
SIAM J. Sci. Comput., 33, 2011.

C. Janna, M. Ferronato, F. Sartoretto, and G. Gambolati. FSAIPACK: A software package
for high-performance factored sparse approximate inverse preconditioning. {\em ACM
Trans. Math. Softw.}, 41(2):10:1{10:26, February 2015.


\end{document}
